\documentclass{article}
\usepackage{xcolor}

\begin{document}

Dado que la población crece según la \textbf{ley logística}, usamos la ecuación diferencial:

\[
\frac{dP}{dt} = a P \left( 1 - \frac{b}{a} P \right)
\]

de aquí se asocia que:\\
- \( \textcolor{blue}{a} \) es la \textbf{tasa de crecimiento intrínseca}.\\
- \( \textcolor{red}{b} \) está relacionado con la capacidad de carga \( K = \frac{a}{b} \).

Los datos que tenemos son:\\
- La población límite es \( K = 5 \times 10^8 \), por lo que \( K = \frac{a}{b} \).\\
- La población se \textbf{duplica cada 40 minutos} cuando es baja.\\
- Queremos encontrar \( P(120) \) (después de 2 horas) para:\\
  - Caso (a): \( P_0 = 10^8 \)\\
  - Caso (b): \( P_0 = 10^9 \)

\textcolor{purple}{\textbf{Paso 1: Determinar \( a \)}}  
Cuando la población es pequeña, sigue un crecimiento exponencial. \qquad \textit{why?... (nótese que para valores muy pequeño el término $\frac{b}{a}$ es despreciable, por lo cual nos queda una ecuación bastante simple de resolver por variables separables, esto nos da qué...}

\[
P(t) \approx P_0 e^{a t}
\]

Como se duplica cada 40 minutos:

\[
2 P_0 = P_0 e^{40a}
\]

Dividiendo por \( P_0 \):

\[
2 = e^{40a}
\]

Aplicando logaritmo natural:

\[
\ln 2 = 40a
\]

\[
a = \frac{\ln 2}{40} \approx \frac{0.693}{40} \approx 0.017325 \, \text{min}^{-1}
\]

\textcolor{purple}{\textbf{Paso 2: Determinar \( b \)}}  
Sabemos que:

\[
K = \frac{a}{b} = 5 \times 10^8
\]

Despejando \( b \):

\[
b = \frac{a}{K} = \frac{0.017325}{5 \times 10^8} = 3.465 \times 10^{-11}
\]

\textcolor{purple}{\textbf{Paso 3: Solución de la ecuación logística}}  
La solución general de:

\[
\frac{dP}{dt} = a P \left( 1 - \frac{b}{a} P \right)
\]

es:

\[
P(t) = \frac{K}{1 + A e^{-a t}}
\]

donde \( A \) se obtiene de la condición inicial \( P(0) = P_0 \):

\[
A = \frac{K}{P_0} - 1
\]

\textcolor{purple}{\textbf{Paso 4: Evaluar para \( t = 120 \) minutos}}  

\textbf{Caso (a): \( P_0 = 10^8 \)}  

\[
A = \frac{5 \times 10^8}{10^8} - 1 = 4
\]

Entonces:

\[
P(120) = \frac{5 \times 10^8}{1 + 4 e^{-0.017325 \times 120}}
\]

Calculamos:

\[
e^{-2.079} \approx 0.125
\]

\[
P(120) = \frac{5 \times 10^8}{1 + 4(0.125)}
\]

\[
= \frac{5 \times 10^8}{1.5} \approx 3.33 \times 10^8
\]

\textbf{Caso (b): \( P_0 = 10^9 \)}  

\[
A = \frac{5 \times 10^8}{10^9} - 1 = -0.5
\]

\[
P(120) = \frac{5 \times 10^8}{1 - 0.5 e^{-2.079}}
\]

\[
P(120) = \frac{5 \times 10^8}{1 - 0.0625}
\]

\[
= \frac{5 \times 10^8}{0.9375} \approx 5.33 \times 10^8
\]

\textcolor{orange}{\textbf{Nótese que}}\\  
- Si \( P_0 = 10^8 \), la población después de 2 horas será aproximadamente \( \mathbf{3.33 \times 10^8} \).\\
- Si \( P_0 = 10^9 \), la población después de 2 horas será aproximadamente \( \mathbf{5.33 \times 10^8} \), alcanzando casi la capacidad de carga.

\end{document}
